\documentclass[11pt]{article}

\usepackage[margin=1in]{geometry}
\usepackage{amsmath,amssymb,amsthm,mathtools}
\usepackage{graphicx}
\usepackage{microtype}
\usepackage[hidelinks]{hyperref}

\newtheorem{theorem}{Theorem}
\newtheorem{lemma}{Lemma}
\newcommand{\R}{\mathbb{R}}
\newcommand{\ellTwo}{\ell_2}
\newcommand{\Lip}{\operatorname{Lip}}
\newcommand{\lipspace}{\operatorname{lip}}
\newcommand{\norm}[1]{\left\lVert #1\right\rVert}

\title{A Reflexive-Subspace Counterexample\\
to Question 5.1 of arXiv:1606.05999}
\author{Run \texttt{fa\_banach\_001}, agent \texttt{agent\_lane\_13}}
\date{22 July 2026}

\begin{document}
\maketitle

\begin{center}
\fbox{\parbox{0.9\textwidth}{
\textbf{Status.} Candidate full negative answer, likely valid, pending human
review. The example does not merely obstruct distortion $1+\varepsilon$:
the resulting little-Lipschitz space admits no isomorphic embedding into
$c_0$ at all.
}}
\end{center}

\section*{Source Question}

García-Lirola, Petitjean, and Rueda Zoca ask the following in Question~5.1
of \cite{GPRZ}, on source PDF page~20:

\begin{quote}
Let $M$ and $\tau$ be under the hypothesis of Theorem 2.9. Is
$lip_\tau(M)$ $(1+\varepsilon)$-isometric to a subspace of $c_0$ for every
$\varepsilon>0$?
\end{quote}

\begin{center}
\includegraphics[width=\textwidth]{figures/open_problem_crop.png}

\small Source PDF page 20, including the complete statement of Question 5.1.
\end{center}

The paper works over the real field. Its standing setup for Theorem~2.9 is a
bounded pointed metric space $(M,d)$ with a compact topology $\tau$ for which
$d$ is $\tau$-lower semicontinuous. The theorem additionally assumes that
$M$ is separable and complete, that $\tau$ is metrizable, and that
\begin{equation*}\tag{P}
 \forall u,v\in M\ \forall \eta>0\ \exists f\in B_{lip_\tau(M)}:
 |f(u)-f(v)|\ge d(u,v)-\eta.
\end{equation*}
Here $lip_\tau(M)$ consists of the pointed $\tau$-continuous Lipschitz
functions satisfying
\[
 \lim_{r\downarrow0}\ \sup_{0<d(u,v)<r}
 \frac{|f(u)-f(v)|}{d(u,v)}=0.
\]

\section*{Proof Intuition}

Snowflaking a norm by the square root turns every bounded linear functional
into a little-Lipschitz function: a norm-linear increment is quadratic in the
new distance. On the weak unit ball of $\ell_2$, all these linear functionals
are also weakly continuous. Consequently $lip_\tau(M)$ contains a closed copy
of the whole reflexive space $\ell_2$.

This is incompatible with an embedding into $c_0$. One standard way to see
the obstruction is that $c_0$ has the hereditary Dunford--Pettis property.
If a closed copy of $\ell_2$ lived in $c_0$, that copy would have the
Dunford--Pettis property; but its identity is weakly compact by reflexivity
and is not completely continuous. Thus the source question fails in a much
stronger form than requested.

\section*{The Counterexample}

\begin{theorem}\label{thm:main}
There is a pair $(M,\tau)$ satisfying all hypotheses of Theorem~2.9 in
\cite{GPRZ} such that $lip_\tau(M)$ is not isomorphic to a subspace of $c_0$.
In particular, Question~5.1 has a negative answer.
\end{theorem}

\begin{proof}
Let $H=\ellTwo$ over $\R$, put $M=B_H$, and use $0$ as the base point. Let
$\tau$ be the weak topology on $B_H$ and define
\[
 d(u,v)=\norm{u-v}_2^{1/2}\qquad(u,v\in B_H).
\]

We first audit the hypotheses. The function $\omega(t)=\sqrt t$ is a
nontrivial gauge in the terminology of \cite{GPRZ}: it is continuous,
increasing, and subadditive, it vanishes at zero, it satisfies
$\omega(t)\ge t$ for $0\le t\le1$, and
$\omega(t)/t\to\infty$ as $t\downarrow0$. The sentence introducing the
paper's corollary immediately after Theorem~2.9 explicitly identifies the
following as a special case:
\[
 \bigl(B_{X^*},\ \omega\circ\norm{\mathord\cdot}\bigr)
\]
with its weak-star topology. Taking $X=H$ and identifying $X^*=H$ places our
example in the family highlighted by the authors. We nevertheless verify
property~(P) directly below, so the counterexample does not depend on that
identification.

For completeness, the remaining topological and metric hypotheses are also
immediate. The ball $B_H$ is weakly compact because $H$ is reflexive, and its
weak topology is metrizable because $H^*=H$ is separable. The snowflaked
metric is bounded and separable. It is complete because a $d$-Cauchy sequence
is norm Cauchy. Finally, the norm is weakly lower semicontinuous, so
$(u,v)\mapsto d(u,v)$ is lower semicontinuous for the product weak topology.

We now prove property~(P). Fix $u,v\in B_H$ with $u\ne v$, put
$\delta=\norm{u-v}_2$, and choose $h\in S_H$ so that
$\langle u-v,h\rangle=\delta$. For $\rho>0$ define
\[
 \phi_\rho(s)=\sqrt{s_++\rho}-\sqrt\rho,
 \qquad s_+=\max\{s,0\}.
\]
The function $\phi_\rho$ is globally Lipschitz in the ordinary metric, with
constant at most $1/(2\sqrt\rho)$. It also satisfies
\begin{equation}\label{eq:holder-ramp}
 |\phi_\rho(s)-\phi_\rho(t)|\le\sqrt{|s-t|}
 \qquad(s,t\in\R).
\end{equation}
Indeed, this is the usual square-root estimate if $s,t\ge0$, it is trivial if
$s,t\le0$, and if $s\le0\le t$, then
\[
 \phi_\rho(t)=\sqrt{t+\rho}-\sqrt\rho
 \le\sqrt t\le\sqrt{|t-s|}.
\]
Set
\[
 f_\rho(z)=\phi_\rho(\langle z-v,h\rangle)
            -\phi_\rho(\langle -v,h\rangle),
 \qquad z\in B_H.
\]
Then $f_\rho(0)=0$, and $f_\rho$ is weakly continuous. From
\eqref{eq:holder-ramp},
\[
 |f_\rho(z)-f_\rho(w)|
 \le\sqrt{|\langle z-w,h\rangle|}
 \le\sqrt{\norm{z-w}_2}=d(z,w),
\]
so $\norm{f_\rho}_{Lip}\le1$. Its ordinary Lipschitz bound gives
\[
 \frac{|f_\rho(z)-f_\rho(w)|}{d(z,w)}
 \le \frac{1}{2\sqrt\rho}\,d(z,w),
\]
which tends uniformly to zero with $d(z,w)$. Hence
$f_\rho\in B_{lip_\tau(M)}$. Finally,
\[
 f_\rho(u)-f_\rho(v)
 =\sqrt{\delta+\rho}-\sqrt\rho
 \longrightarrow\sqrt\delta=d(u,v)
 \qquad(\rho\downarrow0).
\]
This proves property~(P); the case $u=v$ is immediate.

For $x\in H$, define $Jx:M\to\R$ by
\[
 (Jx)(u)=\langle u,x\rangle.
\]
This function vanishes at the base point and is weakly continuous. Moreover,
for $u\ne v$,
\begin{equation}\label{eq:key}
\begin{aligned}
 \frac{|(Jx)(u)-(Jx)(v)|}{d(u,v)}
 &\le \frac{\norm{x}_2\norm{u-v}_2}{\norm{u-v}_2^{1/2}}\\
 &=\norm{x}_2\norm{u-v}_2^{1/2}
 =\norm{x}_2d(u,v).
\end{aligned}
\end{equation}
Because $d(u,v)\le\sqrt2$ on $B_H$, equation \eqref{eq:key} proves that
$Jx$ is Lipschitz and
\[
 \norm{Jx}_{lip_\tau(M)}\le\sqrt2\norm{x}_2.
\]
It also proves the little-Lipschitz condition: if $0<d(u,v)<r$, then the
quotient in \eqref{eq:key} is at most $r\norm{x}_2$, uniformly in $u,v$.
Thus $Jx\in lip_\tau(M)$.

The reverse norm estimate is just as direct. For $x\ne0$, take
$u=x/\norm{x}_2$ and $v=0$. Then $u\in B_H$, $d(u,0)=1$, and
\[
 \frac{|(Jx)(u)-(Jx)(0)|}{d(u,0)}=\norm{x}_2.
\]
Consequently
\begin{equation}\label{eq:embedding}
 \norm{x}_2\le\norm{Jx}_{lip_\tau(M)}
 \le\sqrt2\norm{x}_2 \qquad(x\in H).
\end{equation}
The linear map $J:H\to lip_\tau(M)$ is therefore bounded below, so its range
is a closed subspace isomorphic to $\ell_2$.

Suppose toward a contradiction that $lip_\tau(M)$ admitted an isomorphic
embedding $T$ into $c_0$. Then $T(JH)$ would be a closed subspace of $c_0$
isomorphic to $\ell_2$. But $c_0$ has the hereditary Dunford--Pettis property
\cite{Cembranos}; hence $T(JH)$ has the Dunford--Pettis property. Since it is
reflexive, its identity operator is weakly compact and therefore, by the
Dunford--Pettis property, completely continuous. This is impossible: the
image under an isomorphism of the orthonormal basis of $\ell_2$ is weakly
null while its norms stay bounded away from zero. This contradiction proves
that no such $T$ exists.
\end{proof}

\section*{Verification Notes}

\begin{itemize}
\item \textbf{Property~(P) is explicit.} The functions $f_\rho$ have
Lipschitz norm at most one, are weakly continuous and little Lipschitz, and
their increments converge to the required snowflake distance. The source's
own weak-star-ball corollary independently confirms that this is the intended
class of examples.
\item \textbf{The snowflake is essential.} In the new metric,
$\norm{u-v}_2=d(u,v)^2$, so the Lipschitz quotient of a linear functional is
bounded by $\norm{x}_2d(u,v)$ and tends uniformly to zero at small scales.
\item \textbf{The copy is closed.} The lower inequality in
\eqref{eq:embedding}, not merely injectivity, guarantees that $JH$ is a closed
subspace and permits restriction of any hypothetical embedding into $c_0$.
\item \textbf{The $c_0$ contradiction is isomorphic.} Hereditary
Dunford--Pettis is inherited by every closed subspace of $c_0$; reflexivity
and the Dunford--Pettis property are invariant under isomorphism. Thus no
choice of distortion can evade the obstruction.
\item \textbf{No computational dependency.} Every estimate is exact. The
argument uses no numerical experiment, asymptotic approximation, or
unverified interchange of limits.
\end{itemize}

\section*{Novelty Check and Limitations}

A bounded search through 22 July 2026 covered the run registry, solution,
attempt, and proof-gap indexes; the local source corpus; the arXiv id, title,
and authors; exact wording from Question~5.1; and searches combining
\texttt{lip\_tau(M)}, $c_0$, and the paper's weak-star snowflake notation.
No later answer to Question~5.1 was found. This supports, but does not prove,
novelty.

The result addresses Question~5.1 only. It makes no claim about Question~5.2
in the same paper, concerning the Schur property of a vector-valued
Lipschitz-free space. The recommended human check is the smoothed-ramp proof
of property~(P), especially \eqref{eq:holder-ramp}; the rest of the proof is
elementary Banach-space geometry.

\begin{thebibliography}{9}
\bibitem{GPRZ}
L. García-Lirola, C. Petitjean, and A. Rueda Zoca,
\emph{On the structure of spaces of vector-valued Lipschitz functions},
Rev. R. Acad. Cienc. Exactas Fís. Nat. Ser. A Mat. \textbf{112} (2018),
727--749; \href{https://arxiv.org/abs/1606.05999}{arXiv:1606.05999}.

\bibitem{Cembranos}
P. Cembranos,
\emph{The hereditary Dunford--Pettis property on $C(K,E)$},
Illinois J. Math. \textbf{31} (1987), no.~3, 365--373.
\end{thebibliography}

\end{document}
